Since the power radiated as gravitational waves is the negative of the time rate of
change of the fluid pulsion energy, the \emph{radiated power as measured far from the star is}
%
\begin{equation}
  P = -2\,\tau_{n}^{-1} E_{n,\text{puls}}^{-k/\tau_{n}} \;.
  \label{eq:30}
\end{equation}

\subsection*{b) The Region near the Wave Front}

Thus far we have described mathematically the quasi-normal pulsations only in the
region far behind the wave front. Near the wavefront the description is more complicated because one must superimpose a large number of normal modes with various
amounts of outgoing and incoming radiation in order to obtain a sharp ``turning-on'' of
the perturbation at the wave front.

We shall confine our analysis near the wavefront to those quasi-normal modes with
$|f/\tau_{n}| \ll \sigma_{n}$ ($\sigma_{n}$ near the positive real axis); and we shall pattern our analysis after the
classic analysis of nuclear particle decay by Gamow (1931) and by Breit and Wert
(1935).

We build the wave front and associated quasi-normal pulsations out of a linear
combination of real standing-wave normal modes with eigenfrequencies near $\sigma_{n}$---and
hence also near $\omega_{n} = \sigma_{n} + i/\tau_{n}$. Each of these standing-wave modes is normalized to
have the same real amplitude, $B_{n}$, at the center of the star (cf.\ eq.\ [18]). In the wave
zone reality of the eigenfunctions guarantees that the ingoing and outgoing amplitudes
are complex conjugates of each other (cf.\ eq.\ [18]):
%
\begin{equation}
  C^{(\mathrm{I})} = C^{(\mathrm{O})*} \;.
  \label{eq:31}
\end{equation}

These wave-zone amplitudes are conveniently expressed as functions of the real frequency, $\omega$, by means of a power series expansion in the complex frequency plane about
the purely outgoing frequency $\omega_{n} = \sigma_{n} + i/\tau_{n}$:
%
\begin{equation}
  C_{\omega}^{(\mathrm{O})} = C_{\Omega}^{(\mathrm{O})} = \int dC^{(\mathrm{O})}/d\omega\big|_{\Omega}\,(\omega - \omega_{n} - i/\tau_{n}) \;.
  \label{eq:32}
\end{equation}

The particular combination of real, standing-wave normal modes which yields the desired wave front and subsequent quasi-normal pulsations is this:
%
\begin{equation}
  K(r,t) = \int_{-\infty}^{+\infty}
    \frac{K_{n}(r)}{(\omega - \sigma_{n})^{2} + \tau_{n}^{-2}}
    \frac{1}{2\pi}\,e^{i\omega t}\,d\omega
  + \begin{pmatrix} \text{complex} \\ \text{conjugate} \end{pmatrix} ,
  \label{eq:33}
\end{equation}
%
and similarly for $H_{0}, W, V$.

In the wave zone this becomes, for times\footnote{For times $t < 0$ (eq.\ (30))---incoming gravitational waves
with a trailing wave front. We choose to assume $t = 0$. The perturbation is first turned on inside the star
at time $t = 0$, thereby discarding the incoming solution for $t < 0$.} $t > 0$ [combine eqs.\ (18), (32), and (33); integrate $\omega \to \infty$ since the resonant denominator in eq.\ (33) is huge for negative
$\omega$; close the integration contours at complex infinity; and use the method of residues to
evaluate the integrals]:
%
\begin{equation}
  K(r,t)
  = \begin{cases}
      0 & \text{if} \quad t < r + 2M\ln r \;, \\[4pt]
      \dfrac{1}{2}\,dC^{(\mathrm{O})}/d\omega\big|_{\Omega}\,
        \exp\!\bigl[-(t - r - 2M\ln r)/\tau_{n} + i\sigma_{n}(t - r - 2M\ln r) + \delta\bigr]
      & \text{if} \quad t > r + 2M\ln r \;.
    \end{cases}
  \tag{34a}
\end{equation}
